% ============================================================
%  MR2020-06  1U nut block (proposal, standalone sheet)
%  Printed spacer that runs inside the front slot of the front uprights
%  and holds the three M5 drop-in nuts of one rack unit.
%  pdflatex bloco-1u.tex (twice)
% ============================================================
\documentclass[10pt]{article}
\usepackage[a4paper,margin=0mm]{geometry}
\usepackage[utf8]{inputenc}
\usepackage[T1]{fontenc}
\usepackage{iso}
% ============================================================
%  MINI RACK 2020, cotas mestras (mm)
%  Fonte única de verdade. Alterar aqui recompila tudo.
% ============================================================

% ---------- estrutura ----------
\def\Perfil{20}            % extrusão 2020
\def\RackL{206}            % largura externa
\def\RackP{350}            % profundidade externa
\def\RackA{586}            % altura total da estrutura, conectores incluídos
\def\Conector{3.5}          % acréscimo de cada conector de canto: vale no topo E na base
\def\Coluna{579}           % corte do upright = RackA - 2 x Conector

\def\CantL{20}            % perna da cantoneira L (catálogo), largura 18 sobre a face de 20

% ---------- cortes ----------
\def\TravL{166}            % width rail, 4 peças (base e topo)
\def\TravP{310}            % depth rail, 4 peças (base e topo)

% ---------- rack 8" custom ----------
\def\Furos{186}            % eixo a eixo dos furos
\def\Painel{203.2}         % largura do painel (8 polegadas)
\def\Vao{166}              % vão livre entre colunas

% ---------- unidades ----------
\def\Uu{44.45}             % 1U
\def\NU{12}                % número de unidades
\def\ZonaU{533.4}          % 12U

% Referências verticais
\def\UTopoEstr{26.3}       % topo do U1, a partir do topo da estrutura
\def\UTopoCol{22.8}        % topo do U1, a partir da ponta superior do upright

% furação EIA-310 dentro de 1U (a partir do topo do U)
\def\FuroA{6.35}
\def\FuroB{22.225}
\def\FuroC{38.1}

\pagestyle{empty}

% ---------- block dimensions (mm) ----------
\def\BL{44.45}    % length = 1U
\def\BH{3.9}      % height: 1.4 below the face to 5.3, 0.2 clear of lips and core
\def\BP{6.4}      % nut pocket along the slot: nut 6 + 0.4
\def\BD{2.6}      % pocket depth from the top of the block
\def\BF{1.3}      % floor under the pockets
\def\BT{5.2}      % half width at the top (10.4)
\def\BB{3.7}      % half width at the bottom (7.4)
\def\BM{4.2}      % half width of the floor at the pocket bottom (8.4)

\begin{document}
\begin{folha}

\def\material{PETG, 3D printed, 100\,\% infill}
\def\data{2026-10-29}

\node[txtvista, anchor=west] at (24,280) {1U NUT BLOCK \quad PROPOSAL FOR APPROVAL};

% ============================================================
%  FRONT VIEW 2:1, pockets up (towards the slot opening)
%  sheet point = (FVx + 2 s, FVy + 2 h)
% ============================================================
\def\FVx{62} \def\FVy{250}
\begin{scope}[shift={(\FVx,\FVy)}, x=2mm, y=2mm]
  \draw[linhaA] (0,0) -- (\BL,0) -- (\BL,\BH) -- (41.3,\BH) -- (41.3,\BF) -- (34.9,\BF) -- (34.9,\BH)
    -- (25.425,\BH) -- (25.425,\BF) -- (19.025,\BF) -- (19.025,\BH) -- (9.55,\BH) -- (9.55,\BF)
    -- (3.15,\BF) -- (3.15,\BH) -- (0,\BH) -- cycle;
  \foreach \c in {6.35,22.225,38.1}{\draw[linhaG] (\c,-1.2) -- (\c,{\BH+1.2});}
\end{scope}
% positions of the nut axes from the block end
\cotaHor{\FVx}{\FVx+2*6.35}{\FVy+13}{\FVy+7.8}{\FVy+10.2}{6.35}
\cotaHor{\FVx}{\FVx+2*22.225}{\FVy+19}{\FVy+7.8}{\FVy+10.2}{22.225}
\cotaHor{\FVx}{\FVx+2*38.1}{\FVy+25}{\FVy+7.8}{\FVy+10.2}{38.1}
\cotaHor{\FVx}{\FVx+2*\BL}{\FVy-7}{\FVy}{\FVy}{44.45 (1U)}
\cotaVer{\FVy}{\FVy+2*\BH}{\FVx-7}{\FVx}{\FVx}{3.9}
\cotaVer{\FVy+2*\BF}{\FVy+2*\BH}{\FVx+2*\BL+7}{\FVx+2*41.3}{\FVx+2*\BL}{2.6}
\cotaVer{\FVy}{\FVy+2*\BF}{\FVx+2*\BL+15}{\FVx+2*\BL}{\FVx+2*41.3}{1.3}
\node[txtvista] at (\FVx+44.45,238) {FRONT VIEW \quad 2:1};

% section planes on the front view
\draw[linhaB] ({\FVx+2*14.3},{\FVy-3}) -- ({\FVx+2*14.3},{\FVy-1});
\draw[linhaB] ({\FVx+2*14.3},{\FVy+2*\BH+1}) -- ({\FVx+2*14.3},{\FVy+2*\BH+3});
\node[font=\fontsize{2.6mm}{3mm}\selectfont\bfseries] at ({\FVx+2*14.3+2.5},{\FVy-3}) {A};
\draw[linhaB] ({\FVx+2*22.225},{\FVy-3}) -- ({\FVx+2*22.225},{\FVy-1});
\node[font=\fontsize{2.6mm}{3mm}\selectfont\bfseries] at ({\FVx+2*22.225+2.5},{\FVy-3}) {B};

% ============================================================
%  TOP VIEW 2:1 (first angle: below the front view)
%  sheet point = (TVx + 2 s, TVc + 2 w)
% ============================================================
\def\TVx{62} \def\TVc{212}
\begin{scope}[shift={(\TVx,\TVc)}, x=2mm, y=2mm]
  % silhouette: full top width between pockets, floor width inside them
  \draw[linhaA] (0,-\BT) -- (3.15,-\BT) -- (3.15,-\BM) -- (9.55,-\BM) -- (9.55,-\BT) -- (19.025,-\BT)
    -- (19.025,-\BM) -- (25.425,-\BM) -- (25.425,-\BT) -- (34.9,-\BT) -- (34.9,-\BM) -- (41.3,-\BM)
    -- (41.3,-\BT) -- (\BL,-\BT) -- (\BL,\BT) -- (41.3,\BT) -- (41.3,\BM) -- (34.9,\BM) -- (34.9,\BT)
    -- (25.425,\BT) -- (25.425,\BM) -- (19.025,\BM) -- (19.025,\BT) -- (9.55,\BT) -- (9.55,\BM)
    -- (3.15,\BM) -- (3.15,\BT) -- (0,\BT) -- cycle;
  % pocket walls across the block
  \foreach \s in {3.15,9.55,19.025,25.425,34.9,41.3}{\draw[linhaA] (\s,-\BM) -- (\s,\BM);}
  % bottom edges, hidden under the sloped sides
  \draw[linhaF] (0,-\BB) -- (\BL,-\BB);
  \draw[linhaF] (0,\BB) -- (\BL,\BB);
  % axes
  \draw[linhaG] (-1.5,0) -- ({\BL+1.5},0);
  \foreach \c in {6.35,22.225,38.1}{\draw[linhaG] (\c,{-\BT-1.2}) -- (\c,{\BT+1.2});}
\end{scope}
\cotaHor{\TVx}{\TVx+2*3.15}{\TVc-17}{\TVc-10.4}{\TVc-10.4}{3.15}
\cotaHor{\TVx+2*3.15}{\TVx+2*9.55}{\TVc-17}{\TVc-8.4}{\TVc-8.4}{6.4}
\cotaHor{\TVx+2*9.55}{\TVx+2*19.025}{\TVc-17}{\TVc-10.4}{\TVc-10.4}{9.475}
\cotaVer{\TVc-2*\BT}{\TVc+2*\BT}{\TVx+2*\BL+7}{\TVx+2*\BL}{\TVx+2*\BL}{10.4}
\cotaVer{\TVc-2*\BM}{\TVc+2*\BM}{\TVx+2*\BL+15}{\TVx+2*41.3}{\TVx+2*41.3}{8.4}
\node[txtvista] at (\TVx+44.45,188) {TOP VIEW \quad 2:1};

% ============================================================
%  SECTIONS 4:1 across the slot, block inside the profile cavity
%  profile drawn as adjacent part (phantom); u across, d depth from the face
%  sheet point = (Sx + 4 u, Sy - 4 d)
% ============================================================
\newcommand{\cavidade}{%
  % face, lips (1.2 thick, opening 6.2) and the trapezoid cavity down to the 8 mm core
  \draw[linhaK] (-8,0) -- (-3.1,0) -- (-3.1,-1.2) -- (-5.5,-1.2) -- (-4,-5.5) -- (4,-5.5) -- (5.5,-1.2)
    -- (3.1,-1.2) -- (3.1,0) -- (8,0);
  \draw[linhaG] (0,1) -- (0,-6.5);}

% A-A: between two nuts, full block section
\def\SAx{64} \def\SAy{170}
\begin{scope}[shift={(\SAx,\SAy)}, x=4mm, y=4mm]
  \cavidade
  \filldraw[linhaA, fill=white] (-\BT,-1.4) -- (\BT,-1.4) -- (\BB,-5.3) -- (-\BB,-5.3) -- cycle;
  \path[pattern=north east lines] (-\BT,-1.4) -- (\BT,-1.4) -- (\BB,-5.3) -- (-\BB,-5.3) -- cycle;
\end{scope}
\cotaHor{\SAx-4*\BT}{\SAx+4*\BT}{\SAy+5}{\SAy-5.6}{\SAy-5.6}{10.4}
\cotaHor{\SAx-4*\BB}{\SAx+4*\BB}{\SAy-26.4}{\SAy-21.2}{\SAy-21.2}{7.4}
\cotaVer{\SAy-21.2}{\SAy-5.6}{\SAx+26.4}{\SAx+4*\BB}{\SAx+4*\BT}{3.9}
\cotaVer{\SAy-5.6}{\SAy}{\SAx-26.4}{\SAx-4*\BT}{\SAx-4*3.1}{1.4}
\node[txtvista] at (\SAx,136) {SECTION A-A \quad 4:1};
\node[txtnota] at (\SAx,132) {between two nuts};

% B-B: through a nut, the floor carries the block under the nut
\def\SBx{150} \def\SBy{170}
\begin{scope}[shift={(\SBx,\SBy)}, x=4mm, y=4mm]
  \cavidade
  \filldraw[linhaA, fill=white] (-\BM,-4) -- (\BM,-4) -- (\BB,-5.3) -- (-\BB,-5.3) -- cycle;
  \path[pattern=north east lines] (-\BM,-4) -- (\BM,-4) -- (\BB,-5.3) -- (-\BB,-5.3) -- cycle;
  % drop-in nut, turned across the slot: body 10.5 under the lips, neck 6 in the opening
  \filldraw[linhaA, fill=gray!30] (-3,-0.2) -- (3,-0.2) -- (3,-1.2) -- (5.25,-1.2) -- (4.25,-3.5)
    -- (-4.25,-3.5) -- (-5.25,-1.2) -- (-3,-1.2) -- cycle;
\end{scope}
\cotaHor{\SBx-4*5.25}{\SBx+4*5.25}{\SBy+5}{\SBy-5.6}{\SBy-5.6}{10.5 nut}
\cotaHor{\SBx-4*\BM}{\SBx+4*\BM}{\SBy-26.4}{\SBy-16}{\SBy-16}{8.4}
\cotaVer{\SBy-21.2}{\SBy-16}{\SBx+26.4}{\SBx+4*\BB}{\SBx+4*\BM}{1.3}
\cotaVer{\SBy-16}{\SBy-14}{\SBx-26.4}{\SBx-4*\BM}{\SBx-4*4.25}{0.5}
\node[txtvista] at (\SBx,136) {SECTION B-B \quad 4:1};
\node[txtnota] at (\SBx,132) {through a nut};

% ============================================================
%  STACK IN THE SLOT 1.5:1, top of an upright, section along the slot axis
%  sheet point = (KX + 1.5 s, KY - 1.5 d), s from the upright end
% ============================================================
\def\KX{30} \def\KY{102}
\begin{scope}[shift={(\KX,\KY)}, x=1.5mm, y=1.5mm]
  % core of the profile under the slot, and the face plane
  \fill[pattern=north east lines] (0,-5.5) rectangle (100,-7);
  \draw[linhaK] (0,-5.5) -- (100,-5.5);
  \draw[linhaB] (0,0) -- (100,0);
  \draw[linhaA] (0,1.5) -- (0,-7);
  % corner connector arm in the first 20 mm
  \filldraw[linhaA, fill=conector] (0,-0.5) rectangle (20,-5.3);
  % block U1 from 22.8, block U2 from 67.25 broken at 100
  \begin{scope}
    \clip (0,-8) rectangle (100,3);
    \foreach \o in {22.8,67.25}{
      \begin{scope}[shift={(\o,0)}]
        \draw[linhaA] (0,-5.3) -- (0,-1.4) -- (3.15,-1.4) -- (3.15,-4) -- (9.55,-4) -- (9.55,-1.4)
          -- (19.025,-1.4) -- (19.025,-4) -- (25.425,-4) -- (25.425,-1.4) -- (34.9,-1.4) -- (34.9,-4)
          -- (41.3,-4) -- (41.3,-1.4) -- (44.45,-1.4) -- (44.45,-5.3) -- cycle;
        \foreach \c in {6.35,22.225,38.1}{
          \filldraw[linhaA, fill=gray!30] ({\c-3},-0.2) rectangle ({\c+3},-3.5);
          \draw[linhaG] (\c,1.5) -- (\c,-7);}
      \end{scope}}
  \end{scope}
\end{scope}
\quebraV{\KX+150}{\KY-11}{\KY+2}
\cotaHor{\KX}{\KX+1.5*20}{\KY+8}{\KY+2}{\KY+2}{20}
\cotaHor{\KX}{\KX+1.5*22.8}{\KY+14}{\KY+2}{\KY+2}{22.8}
\cotaHor{\KX}{\KX+1.5*29.15}{\KY+20}{\KY+2}{\KY+2}{29.15}
\cotaHor{\KX+1.5*22.8}{\KX+1.5*67.25}{\KY-16}{\KY-8}{\KY-8}{44.45}
\chamada{\KX+15}{\KY-4.5}{\KX+8}{\KY-24}{north}{CONNECTOR ARM}
\chamada{\KX+1.5*29.15}{\KY-3}{\KX+1.5*29.15+18}{\KY+24}{south}{M5 DROP-IN NUT}
\chamada{\KX+1.5*45}{\KY-5}{\KX+1.5*45}{\KY-24}{north}{BLOCK U1}
\chamada{\KX+1.5*85}{\KY-5}{\KX+1.5*85}{\KY-24}{north}{BLOCK U2}
\node[txtvista, anchor=west] at (128,126) {STACK IN THE FRONT SLOT \quad 1.5:1};

% ============================================================
%  NOTES
% ============================================================
\node[anchor=north west, txtnota, text width=174mm, align=left, inner sep=0pt] at (24,68) {%
  \hyphenpenalty=10000\exhyphenpenalty=10000
  {\fontsize{2.6mm}{3mm}\selectfont\bfseries NOTES}\\[1mm]
  1. All dimensions in mm. Proposal for approval: not modeled yet.\\[0.6mm]
  2. 12 identical blocks per front upright, 24 in all, slid in from the upright end before the corner connector goes on. Print floor down, pockets up, no supports.\\[0.6mm]
  3. Holds the M5 drop-in nut 10.5 × 6 × 3.3 (neck 6 × 6 × 1), turned across the slot: its 6 mm side runs along the slot, in a 6.4 pocket.\\[0.6mm]
  4. Profile cavity assumed 11 wide under the 1.2 lips and 5.5 deep from the face to the core, 0.2 to 0.3 clear on every side: to be confirmed with a test print.\\[0.6mm]
  5. The connector arm takes the first 20 mm of the slot. The stack rests on the lower arm, so the U grid sits 2.8 lower than on sheet 02; the pitch is unchanged.};

\legenda{1U NUT BLOCK}{Printed spacer for the M5 drop-in nuts in the front slot, proposal}{MR2020-06}{2:1, 4:1}{1}{1}

\end{folha}
\end{document}
